\documentclass[11pt,a4paper]{article}

\usepackage[T1]{fontenc}
\usepackage[utf8]{inputenc}
\usepackage[spanish,es-noshorthands]{babel}
\usepackage{amsmath,amssymb}
\usepackage[margin=2.5cm]{geometry}
\usepackage{enumitem}

\title{Los tres balances: guion de derivación}
\author{Mecánica del Continuo}
\date{}

\begin{document}
\maketitle

\noindent\textbf{Receta común a los tres:}
\[
\text{Axioma sobre } V(t)
\ \longrightarrow\ \text{Reynolds}
\ \longrightarrow\ \text{Cauchy}
\ \longrightarrow\ \text{Gauss}
\ \longrightarrow\ V \text{ arbitrario} \Rightarrow \text{localizar}
\]

\vspace{1em}

%======================================================================
\section*{1. Conservación de masa (continuidad)}

Postulo que la masa de un volumen material no cambia:
\[
\frac{D}{Dt}\int_{V(t)}\rho\,dV=0 ,
\]
donde $\rho$ es la \emph{densidad} (masa por unidad de volumen), no la masa.

Aplico el teorema del transporte de Reynolds y me quedan dos efectos:
\[
\int_{V}\Big[\underbrace{\frac{D\rho}{Dt}}_{\substack{\text{cuánto cambia la densidad}\\ \text{siguiendo a la partícula}}}
+\underbrace{\rho\,\frac{\partial v_k}{\partial x_k}}_{\substack{\text{cuánto cambia}\\ \text{su volumen}}}\Big]dV=0 .
\]

Como el volumen $V$ es \textbf{arbitrario}, el integrando debe anularse (localización):
\[
\boxed{\ \frac{D\rho}{Dt}+\rho\,\frac{\partial v_k}{\partial x_k}=0\ }
\]

El resultado dice que densidad y volumen cambian en \textbf{sentidos opuestos}, de forma
que la masa quede igual:
\[
\frac{D\rho}{Dt}=-\rho\,\frac{\partial v_k}{\partial x_k},
\qquad\text{donde}\qquad
\frac{\partial v_k}{\partial x_k}=\nabla\cdot\boldsymbol v=\frac{1}{dV}\frac{D(dV)}{Dt}.
\]

Si además impongo \textbf{incompresibilidad} ($D\rho/Dt=0$, hipótesis adicional):
\[
\boxed{\ \nabla\cdot\boldsymbol v=\frac{\partial v_k}{\partial x_k}=0\ }
\]

\medskip
\noindent\textbf{Cuidado:} $\partial\rho/\partial t=0$ es \emph{estacionario};
$D\rho/Dt=0$ es \emph{incompresible}. Son hipótesis independientes.

%======================================================================
\section*{2. Balance de cantidad de movimiento lineal}

Postulo la segunda ley de Newton: la tasa de cambio de la cantidad de movimiento
de un volumen material iguala a las fuerzas de superficie más las de volumen.
Es un \textbf{balance}, no una conservación (el momento cambia si hay fuerzas).
\[
\frac{D}{Dt}\int_{V(t)}\rho v_i\,dV=\int_{S(t)}t_i\,dS+\int_{V(t)}b_i\,dV
\]

\paragraph{Lado izquierdo.} Aplico Reynolds. Los términos que tienen $v_i$ como
factor común forman la ecuación de continuidad y \textbf{se anulan} (acá uso
conservación de masa):
\[
\frac{D}{Dt}\int_V\rho v_i\,dV
=\int_V\Big[\rho\frac{Dv_i}{Dt}
+v_i\underbrace{\Big(\frac{D\rho}{Dt}+\rho\frac{\partial v_k}{\partial x_k}\Big)}_{=\,0}\Big]dV
=\int_V\rho\frac{Dv_i}{Dt}\,dV
\]

\paragraph{Lado derecho.} La tracción depende de la orientación del plano,
$t_i=t_i(\boldsymbol x,\boldsymbol n)$, así que primero uso el \textbf{teorema de Cauchy}
para introducir el tensor de tensiones,
\[
t_i=\sigma_{ij}n_j ,
\]
y \emph{recién entonces} aplico \textbf{Gauss} para pasar la integral de superficie a volumen:
\[
\int_S t_i\,dS=\int_S\sigma_{ij}n_j\,dS=\int_V\frac{\partial\sigma_{ij}}{\partial x_j}\,dV
\]

\paragraph{Localización.} Como $V$ es arbitrario:
\[
\boxed{\ \rho\,\frac{Dv_i}{Dt}=\frac{\partial\sigma_{ij}}{\partial x_j}+b_i\ }
\qquad\Longleftrightarrow\qquad
\rho\,\dot{\boldsymbol v}=\nabla\cdot\boldsymbol\sigma+\boldsymbol b
\]

\medskip
\noindent El balance de momento \textbf{angular} no agrega una ecuación nueva:
da la restricción $\sigma_{ij}=\sigma_{ji}$.

%======================================================================
\section*{3. Balance de energía (primer principio)}

Postulo el primer principio de la termodinámica: la tasa de cambio de la energía
total (interna $+$ cinética) de un volumen material es igual a la potencia de las
fuerzas de superficie y de volumen, más el calor generado internamente
($q$, \emph{fuente volumétrica}), menos el calor que se escapa por el contorno
($h_i n_i$, con normal saliente, por eso va con menos):
\[
\frac{D}{Dt}\int_{V(t)}\rho\Big(\varepsilon+\tfrac12 v_kv_k\Big)dV
=\underbrace{\int_S t_iv_i\,dS+\int_V b_iv_i\,dV}_{\text{potencia mecánica}}
+\underbrace{\int_V q\,dV-\int_S h_in_i\,dS}_{\text{potencia calórica}}
\]

\paragraph{Lado izquierdo.} Por el lema de conservación de masa, $\rho$ sale afuera
de la derivada:
\[
\frac{D}{Dt}\int_V\rho\,\psi\,dV=\int_V\rho\,\frac{D\psi}{Dt}\,dV
\qquad\text{con}\qquad \psi=\varepsilon+\tfrac12 v_kv_k
\]

\paragraph{Lado derecho.} Uso Cauchy ($t_i=\sigma_{ij}n_j$) y Gauss para pasar las
integrales de superficie a volumen. Como $V$ es arbitrario, localizo:
\[
\rho\frac{D\varepsilon}{Dt}+\rho v_i\frac{Dv_i}{Dt}
=v_i\frac{\partial\sigma_{ij}}{\partial x_j}+\sigma_{ij}\frac{\partial v_i}{\partial x_j}
+b_iv_i+q-\frac{\partial h_i}{\partial x_i}
\]

\paragraph{Resta de la energía mecánica.} Multiplico la ecuación de Cauchy por $v_i$
(teorema de la energía mecánica),
\[
\rho v_i\frac{Dv_i}{Dt}=v_i\frac{\partial\sigma_{ij}}{\partial x_j}+b_iv_i ,
\]
y se la resto: se cancelan la energía cinética y el trabajo de $\sigma$ y $b$, y queda
la ecuación de la \textbf{energía interna}, donde sobrevive $\sigma_{ij}\,\partial v_i/\partial x_j$.

Como $\sigma_{ij}$ es simétrico y la parte antisimétrica del gradiente de velocidad
($W_{ij}$) es la rotación rígida —que no hace trabajo, $\sigma_{ij}W_{ij}=0$— queda
$\sigma_{ij}V_{ij}$:
\[
\boxed{\ \rho\,\frac{D\varepsilon}{Dt}=\sigma_{ij}V_{ij}+q-\frac{\partial h_i}{\partial x_i}\ }
\]

%----------------------------------------------------------------------
\subsection*{Caso particular: conducción en un sólido rígido}

Si el cuerpo es \textbf{rígido} ($V_{ij}=0$), está \textbf{en reposo}
($D/Dt\to\partial/\partial t$) y es \textbf{estacionario} ($\partial/\partial t=0$):
\[
0=q-\frac{\partial h_i}{\partial x_i}
\]

Con la ley de \textbf{Fourier} $h_i=-k\,\partial T/\partial x_i$ y $k$ uniforme,
los dos signos menos se cancelan:
\[
-\frac{\partial h_i}{\partial x_i}
=-\frac{\partial}{\partial x_i}\Big(-k\frac{\partial T}{\partial x_i}\Big)
=+k\frac{\partial^2 T}{\partial x_i\partial x_i}
\]
\[
\boxed{\ k\,\nabla^2 T+q=0\ }
\]

\medskip
\noindent Los dos signos menos son \textbf{independientes}: el de
$-\partial h_i/\partial x_i$ es \emph{geométrico} (normal saliente); el de
$h_i=-k\,\partial_i T$ es \emph{físico} (el calor va de caliente a frío, segunda ley).

%======================================================================
\section*{Resumen}

\begin{center}
\begin{tabular}{@{}llll@{}}
\hline
 & Se balancea & Axioma & Resultado local \\
\hline
Masa & $\rho$ & $\dfrac{D}{Dt}\displaystyle\int\rho\,dV=0$
     & $\dfrac{D\rho}{Dt}+\rho\dfrac{\partial v_k}{\partial x_k}=0$ \\[2ex]
Mom.\ lineal & $\rho v_i$ & Newton
     & $\rho\dfrac{Dv_i}{Dt}=\dfrac{\partial\sigma_{ij}}{\partial x_j}+b_i$ \\[2ex]
Mom.\ angular & $\rho\,\epsilon_{ijk}x_jv_k$ & Newton (momentos)
     & $\sigma_{ij}=\sigma_{ji}$ \\[2ex]
Energía & $\rho(\varepsilon+\tfrac12 v^2)$ & 1.\textsuperscript{er} principio
     & $\rho\dfrac{D\varepsilon}{Dt}=\sigma_{ij}V_{ij}+q-\dfrac{\partial h_i}{\partial x_i}$ \\
\hline
\end{tabular}
\end{center}

\vspace{1em}
\noindent Y después: \textbf{$+$ ley constitutiva $=$ ecuación de campo}
(Navier / Navier--Stokes / calor).

\end{document}
